\scanpage{44}
of ``stratified spaces'' (in a sense which
will be specified below) whose
``\unsure{stones|pieces}'' are $K(\unsure{\pi|\Pi},1)$, \deleted{\unsure{that should}}\inserted{such as}
the Teichmüller moduli spaces
and their compactifications, and \unsure{the|their} subspaces
\inserted{and tubular neighborhoods among these}
that one can associate with them thanks to their
canonical stratifications, the homotopy
type should be expressible
in terms of a suitable inductive
system (most often with a finite index
set) of fundamental groupoids
\unsure{of the said|of the said} \unsure{stones|pieces}.

It seems to me moreover in retrospect,
on reflection, that this theory ``of integration''
of homotopy types, which I nebulously
saw as linked to the great yoga of
relations between homotopy structures (relative
over a topos say) and $n$-groupoids,
is in fact relatively trivial, and
probably amenable to the most standard $\frac12$-
simplicial techniques. I shall
have to come back to this later -
% Pale bottom-left foliation in another hand; no authorial page number at the top.
\margin{43}
